\section{Numerical experiments}
\label{sec:experiments}

Everything above this section is a proof, and proofs leave two things open that
only measurement closes.

The first is whether the code satisfies the identities. A proof establishes them in
exact arithmetic; the implementation reaches them through machinery the derivation
does not mention --- a triangular factor held as packed columns addressed by hand,
orthogonal updates applied by library calls that write in place into a view of
$J$'s own buffer, correct only while that view stays contiguous. A sign error in
the packed indexing, or a rank-one update handed a strided view and quietly given a
copy to modify, would leave every proposition above true and the running solver
wrong. Nothing in the derivation can detect that, and the invariants are exactly
the kind of property that decays quietly: they are maintained across a run rather
than recomputed, so an error does not announce itself but accumulates, and a solver
that has drifted from $J\T G J = I$ still returns a vector that is very nearly
right.

The second is the set of quantities this paper reasons about but cannot predict.
Section~\ref{sec:pdas} argues that guessing the active set is unguaranteed but
sound; how often the guess is certifiable, and how many repairs it takes, are
questions about behaviour. Section~\ref{sec:sweep} derives a cost;
what a warm-started solve costs is not the same claim.

All measurements below are on the implementation~\cite{cvxquadprog} at the version
recorded there, and are produced by scripts distributed with this paper's sources.
Where a claim in an earlier section rests on one of them, it points here.

\subsection{The identities hold in the running code}
\label{ssec:identities}

Every identity stated in this paper is evaluated on both sides against the shipped
implementation, on the internal state that implementation actually holds --- its
$J$, its packed $R$, and the $d$, $z$ and $r$ its own routine returns --- rather
than against a reimplementation written for the purpose. Reaching into the private
interior is the point: the claims are about objects that are internal by design,
and a check written against the public surface could not see them.

Residuals are reported in the scaled sup-norm
\begin{equation}
  \mathrm{rel}(P,Q) = \frac{\norm{P - Q}_\infty}{\max(1, \norm{Q}_\infty)},
  \label{eq:relresid}
\end{equation}
so that figures are comparable across identities of different magnitudes. The
sweep runs over problems of size $n \in \{\qpIdentSizes\}$, several seeds each,
driving the factorisation into every intermediate working set reached by inserting
columns one at a time: \qpIdentStates{} distinct factorisation states over
\qpIdentProblems{} problems, and \qpIdentChecks{} individual checks.

\begin{table}[ht]
\centering
\small
\caption{Worst relative residual~\eqref{eq:relresid} for each of the
\qpIdentCount{} identities, over \qpIdentChecks{} checks on \qpIdentStates{}
factorisation states. Two rows are exactly zero; see the text.}
\label{tab:identities}
\begin{tabular}{llrr}
\toprule
Result & Identity & Worst residual & Checks \\
\midrule
\quadprogIdentityRows
\bottomrule
\end{tabular}
\end{table}

The worst residual anywhere in Table~\ref{tab:identities} is $\qpIdentWorst$. That
is the scale to expect: these are exact-arithmetic identities evaluated in double
precision on matrices of condition number $O(n)$, so a few multiples of the machine
epsilon is agreement, not approximation. No threshold appears anywhere in the
procedure --- the residuals are reported, not tested against a bound.

Two rows are exactly zero, which is a stronger statement than a small residual.

\paragraph{The sign conditions are never violated.} The multipliers of active
inequalities at the returned point are non-negative in every instance, with no
negative excursion at any magnitude. That is what one wants from
Proposition~\ref{prop:sufficient}: dual feasibility is maintained
combinatorially, by the ratio test~\eqref{eq:t1} refusing to step past a zero
crossing, rather than restored afterwards by clipping. A method that enforced it by
clipping would show a small nonzero here.

\paragraph{Rescaling a constraint changes the trajectory not at all.}
Proposition~\ref{prop:scale} claims that the choice made by the selection
rule~\eqref{eq:select} is invariant under $(c_i,b_i) \mapsto (\gamma c_i, \gamma
b_i)$. The measured statement is stronger: over three decades of $\gamma$ either
side of unity, the solver performs an identical number of additions and removals
and returns a minimiser agreeing to $3.3\times10^{-16}$. The invariance is
therefore a property of the entire walk and not only of the selected index, which
is what makes it safe for a caller to write constraints in whatever units are
natural.

\subsection{The guessed active set, and what the certificate catches}
\label{ssec:pdas}

Section~\ref{sec:pdas} makes two claims that only measurement settles: that the
block exchange converges in a number of repairs that does not grow with $n$, and
that the loss of the $P$-matrix property for general $C$ is a real exposure rather
than a technicality. Figure~\ref{fig:pdas} addresses both over
\qpPdasFamilies{} constraint families, \qpPdasInstances{} instances each, at
$n \in \{\qpPdasSizes\}$. The instrumentation wraps the shipped routines rather
than reimplementing them, so what is counted is what runs.

\begin{figure}[ht]
\centering
\includegraphics[width=\textwidth]{graphs/quadprog_pdas.pdf}
\caption{Left: the fraction of guesses the certificate accepts. Right: set repairs
to convergence. Four families are certified without exception at every size and
settle in \qpPdasRepairsLo{} to \qpPdasRepairsHi{} repairs while the exact walk's
outer iterations grow linearly in $n$. The fifth is constructed so that a guessed
set can be linearly dependent, and it fails --- which is the point.}
\label{fig:pdas}
\end{figure}

On the four well-posed families --- bounds, budget-plus-bounds, dense $C$, and
$C$ with equalities --- every guess is certified at every size, and the repair
count rises only from \qpPdasRepairsLo{} to \qpPdasRepairsHi{} across a
sixteen-fold range of $n$, never exceeding \qpPdasRepairsMax{}. The comparison
that matters is against the exact walk on the same instances: on
budget-plus-bounds it takes \qpPdasBudgetOuter{} outer iterations at
$n = \qpPdasRefN{}$, and its count grows linearly where the repair count is very
nearly flat. That is the trade Section~\ref{sec:pdas} describes, measured.

\paragraph{The rank-deficient guess is real, and it is not visible by accident.}
Over \qpPdasInstances{} instances at each of five sizes, not one guess on those
four families produced a rank-deficient working-set system. A reader might
conclude the exposure is theoretical. It is not; it is merely invisible in random
data, because a random dense $C$ almost never has a dependent subset that the
guess happens to name. The fifth family makes it reachable in the most elementary
way, by repeating columns of $C$, so that a guess naming both copies of a
constraint induces a singular system. There the Cholesky rank test fires on
\qpPdasSingularDup\% of attempts, and the certified fraction collapses from
\qpPdasCertSmall\% at $n = \qpPdasNsmall{}$ to \qpPdasCertBig\% at
$n = \qpPdasNbig{}$ --- against \qpPdasSingularBox\%, \qpPdasSingularBudget\%,
\qpPdasSingularDense\% and \qpPdasSingularEq\% on the others. Duplicated
constraints are not exotic; they arise whenever a model is assembled
programmatically from overlapping rules. This is the concrete form of the
structural point in Section~\ref{sec:pdas}: with bound constraints the system is a
principal submatrix of $G^{-1}$ and cannot be singular, and nothing about general
$C$ preserves that.

\paragraph{Where the failures are caught.} Of the candidates that reached the
certificate, \qpPdasAccepted{} were accepted, with a worst KKT
residual~\eqref{eq:certificate} of $\qpPdasAcceptWorst$, and
\qpPdasRejected{} were rejected. We report that asymmetry rather than smoothing
it: on these families the certificate never had to reject a converged candidate,
because the failures were caught earlier and more cheaply, by the Cholesky
factorisation of the working-set system refusing a dependent guess. The
certificate remains the guarantee --- it is what makes the method sound rather
than merely usually right --- but on this evidence it is the rank test that does
the work in practice, and a reader should not infer from
Section~\ref{sec:pdas} that non-optimal converged candidates are common.

\subsection{What a warm-started solve costs}
\label{ssec:sweepcost}

Section~\ref{sec:sweep} derives the cost of recovery from cached factors as
$O(n^2)$, all of it in $x_u = J J\T a$, with $O(nk + k^2)$ for everything after.
Measuring that requires care, because the natural experiment cannot see it: on a
box-constrained family the active set grows with $n$, so $O(nk)$ and $O(n^2)$
predict the same curve and either reading fits. Figure~\ref{fig:sweep} therefore
holds the active set fixed at $k = \qpSweepFixedK{}$ while $n$ grows, where the two
predictions separate --- doubling versus quadrupling per doubling of $n$.

\begin{figure}[ht]
\centering
\includegraphics[width=\textwidth]{graphs/quadprog_sweep.pdf}
\caption{Left: cost per solve against $n$, with the active set held at
$k = \qpSweepFixedK{}$ and with $k$ growing, against a cold solve and two
reference slopes. Right: the local exponent $d\log t / d\log n$, which is what
distinguishes the two regimes.}
\label{fig:sweep}
\end{figure}

The local exponent runs from $\qpSweepSlopeLo$ at the small end to
$\qpSweepSlopeHi$ at the large end. Neither figure is $1$, and the second is the
one the derivation predicts: at $n = \qpSweepNhi{}$ a hit costs
$\qpSweepHitHi\,\mu$s against $\qpSweepHitLo\,\mu$s at
$n = \qpSweepNlo{}$, growing faster than any linear reading allows. So the
$O(n^2)$ of Section~\ref{sec:sweep} is what a large warm solve costs, and an
$O(nk)$ reading of the same formulas --- tempting, because every term after $x_u$
does scale with $k$ --- is wrong.

The small end is worth as much attention, because it is where the class is most
used and where the flat curve invites the wrong explanation. Below roughly
$n = 400$ the exponent is near zero: a hit costs what its dozen array operations
cost to dispatch, not what its arithmetic costs. The flatness is an overhead floor
and it ends where the $n^2$ overtakes dispatch. Reported as a speedup against
solving each member of the family from scratch, reuse is worth
$\qpSweepSpeedLo\times$ at $n = \qpSweepNlo{}$ and $\qpSweepSpeedHi\times$ at
$n = \qpSweepNhi{}$, and the cached set stayed optimal on \qpSweepHitRate\% of
solves under a relative perturbation of $\qpSweepPerturb$ to the linear term.

\subsection{Against solvers that are not active-set methods}
\label{ssec:compare}

The experiments so far compare the solver with its own claims, which establishes
that it is implemented correctly and says nothing about whether the method is
worth using. That is a comparative question, and Table~\ref{tab:compare} and
Figure~\ref{fig:compare} answer it against three alternatives: the reference C
implementation of the same algorithm, an operator-splitting method (OSQP), and an
interior-point method (Clarabel).

A table of times alone would be misleading here, and in a direction that flatters
whichever solver was configured loosest. An active-set method terminates at an
exactly feasible point of the active-set lattice and returns a KKT point to
machine precision; a splitting method stops when a tolerance is met; an
interior-point method approaches the boundary asymptotically and never reaches it.
Every solver is therefore asked for the same accuracy, $\qpCmpTol$, and the
residual it \emph{achieves} is reported beside its time, in the same scaled
sup-norm the certificate~\eqref{eq:certificate} uses. Multipliers are recovered
from the returned iterate by least squares rather than taken from the solver, so
that no solver is charged for reporting them in another convention --- the most
favourable reading of each returned point.

\begin{table}[ht]
\centering
\small
\caption{Time per solve in milliseconds and achieved KKT residual, at
$n \in \{\qpCmpSizes\}$. The C reference row is omitted where no wheel for it is
installed.}
\label{tab:compare}
\begin{tabular}{lrrrrrrrr}
\toprule
& \multicolumn{2}{c}{$n = 50$} & \multicolumn{2}{c}{$n = 100$}
& \multicolumn{2}{c}{$n = 200$} & \multicolumn{2}{c}{$n = 400$} \\
\cmidrule(lr){2-3}\cmidrule(lr){4-5}\cmidrule(lr){6-7}\cmidrule(lr){8-9}
Solver & ms & resid. & ms & resid. & ms & resid. & ms & resid. \\
\midrule
\quadprogCompareRows
\bottomrule
\end{tabular}
\end{table}

\begin{figure}[ht]
\centering
\includegraphics[width=\textwidth]{graphs/quadprog_compare.pdf}
\caption{Time per solve against $n$ on three constraint shapes. The reference
implementation of the same algorithm wins at the small end, where cost is
interpreter dispatch rather than arithmetic, and loses at the large end, where it
is arithmetic.}
\label{fig:compare}
\end{figure}

Three things in the table are worth drawing out.

\paragraph{The same algorithm in two languages crosses over.} The C reference is
the fastest solver here at $n = 50$ and among the slowest at
$n = \qpCmpNhi{}$, where it is $\qpCmpVsRef\times$ slower than the
implementation measured in this paper. Nothing algorithmic separates them: they
walk the same active sets in the same order. What separates them is that below a
few hundred variables a solve costs what its operations cost to dispatch, and
above it a solve costs arithmetic, which reaches vendor-tuned kernels in one and
hand-written scalar loops in the other. This reproduces, independently and as a
by-product, the crossover the implementation is designed around.

\paragraph{Accuracy is where the interior-point method pays.} Clarabel is asked
for $\qpCmpTol$ and delivers between $10^{-9}$ and
$\qpCmpResidClarabel$, against $\qpCmpResidExact$ for the active-set method ---
several orders of magnitude, and not a defect of the solver. The optima here are
vertices of the feasible polyhedron, with many constraints active, and a vertex is
exactly the case an interior-point method approaches asymptotically. OSQP reaches
$\qpCmpResidOsqp$, comparable to the active-set method, but only with its
polishing step enabled and its tolerance set to $\qpCmpTol$; at its default
tolerance it is faster and not comparable. The general point is that on
combinatorially structured optima an active-set method returns an answer the
others can only approach, and comparing times without residuals hides that.

\paragraph{Guessing the active set is worth more than the language.} The
\texttt{fast} path is the quickest of everything measured at
$n = \qpCmpNhi{}$ on every family, and it beats the exact walk by roughly a
factor of two to three --- more than the gap between the two implementations of the
exact walk at the same size. Iterations, not the cost of an iteration, are what
dominate at these sizes, which is the argument of Section~\ref{sec:pdas} arriving
from a different direction.

\paragraph{What this comparison is not.} These are dense problems of small to
moderate size whose optima are vertices, which is the territory
Section~\ref{sec:intro} claims for the method and precisely the territory OSQP and
Clarabel are not built for. Both are designed for large sparse problems, both
exploit sparsity this experiment does not give them, and both would be expected to
win outright at $n$ in the tens of thousands with a sparse $G$. The table is
evidence that the method is the right choice on its own ground, and no evidence at
all about anyone else's.

\subsection{What these experiments do not establish}

Three limits, stated because the section is easy to over-read.

They verify agreement between the code and the identities proved here. They do not
verify the absence of defects elsewhere --- in the constraint-selection scan, the
infeasibility verdict, or the handling of degenerate input. A stronger instrument
for that is a differential comparison of complete active-set trajectories against
an independent implementation of the same method, which is not attempted here.

They say nothing about conditioning. Problems are generated with $G = BB\T + nI$
for Gaussian $B$, positive definite by construction and moderately conditioned;
the residuals in Table~\ref{tab:identities} would grow with $\kappa(G)$, and that
table should not be read as a claim about ill-conditioned instances.

That choice also biases Section~\ref{ssec:pdas} in a direction worth naming,
because it works against the conclusion drawn there. A well-conditioned $G$ has a
minimiser that spreads across coordinates, so few bounds bind and the active set
is small; a realistic covariance, being close to low-rank plus diagonal,
concentrates the optimum and binds many more. The exact walk's iteration count
grows with the size of the set it must reach, so the counts reported here are at
the low end of what an application would produce, and the advantage measured for
the guessed set is correspondingly the low end of the advantage available. The
direction of the bias is favourable to the method this paper does \emph{not}
advocate, which is the safe direction for it to run in, but a reader should not
read \qpPdasBudgetOuter{} outer iterations at $n = \qpPdasRefN{}$ as typical of a
portfolio problem.

Finally, the identity checks are dense and of modest size, chosen for breadth of
shape rather than scale, and the timings are single-machine. An identity that
survives $n = 120$ with a working set of $60$ will not fail at $n = 1200$ for
mathematical reasons; if it fails there it will be for reasons of memory layout,
which the derivation is indifferent to and a timing on one host cannot settle.
